\section{Metodología de Desarrollo e Integración}

El desarrollo de software criptográfico impone requisitos críticos de exactitud matemática y seguridad en la gestión de memoria. Un mínimo error en la manipulación de bits o en la aritmética modular puede resultar en vulnerabilidades catastróficas. Por ello, este proyecto descartó el desarrollo empírico tradicional en favor de una metodología orientada a la verificación continua y la segregación arquitectónica.

\subsection{Arquitectura Orientada a Interfaces}
Para permitir el desarrollo en paralelo y mitigar los bloqueos entre los miembros del equipo, el sistema se diseñó bajo el principio de Segregación de Interfaces (SOLID). Se definieron contratos estrictos (\textit{interfaces} en Java) para los tres componentes principales del estándar: el módulo de \textit{hashing} criptográfico (\texttt{HashService}), la generación de firmas (\texttt{SignatureGenerator}) y la verificación de firmas (\texttt{SignatureVerifier}). 

Esta abstracción garantizó que cada módulo interactuara con los demás únicamente a través de tipos de datos puros y Objetos de Transferencia de Datos (DTOs), eliminando dependencias circulares y facilitando el aislamiento necesario para las pruebas de calidad.

\subsection{Desarrollo Guiado por Pruebas (TDD) con JUnit 5}
La piedra angular de la metodología del proyecto fue la adopción estricta del Desarrollo Guiado por Pruebas (TDD, por sus siglas en inglés). En lugar de programar la lógica del estándar y probarla al final, se codificaron primero las pruebas unitarias utilizando el marco de trabajo JUnit 5. 

El marco de aserción se construyó extrayendo los vectores de prueba oficiales (entradas estáticas y sus salidas matemáticas esperadas) publicados en los apéndices de los documentos GOST R 34.10-2012 y RFC 6986. De esta forma, el código de implementación se consideraba completo única y exclusivamente cuando lograba satisfacer las aserciones estandarizadas.

Para asegurar la legibilidad y mantenibilidad del conjunto de pruebas, cada test unitario se estructuró siguiendo el patrón \textbf{AAA (Arrange, Act, Assert)}:
\begin{itemize}
    \item \textbf{Arrange (Preparar):} Se inicializan los objetos, se configuran las dependencias y se declaran las constantes estáticas (vectores de prueba) extraídas del estándar.
    \item \textbf{Act (Ejecutar):} Se invoca el método principal del componente criptográfico.
    \item \textbf{Assert (Afirmar):} Se evalúa si el resultado de la ejecución coincide exactamente con el vector de control.
\end{itemize}

El siguiente fragmento de código ilustra la aplicación del patrón AAA en la validación del flujo completo del Emisor (Algoritmo I). Para lograr una prueba determinista y replicable, se utilizó una sobrecarga del método que permite inyectar el número aleatorio efímero ($k$), evaluando así el ``camino feliz'' contra los valores exactos del Apéndice A.1:

\begin{lstlisting}[language=Java, caption={Prueba unitaria del Emisor estructurada bajo el patrón AAA (cadenas hexadecimales truncadas por brevedad)}, label={lst:aaa_pattern}]
@Test
public void testSignWithAppendixA1Vectors() {
    // 1. Arrange (Preparar)
    // Configuracion de parametros de dominio (256 bits)
    BigInteger p = new BigInteger("80000000...0431", 16);
    BigInteger a = new BigInteger("7", 16);
    BigInteger b = new BigInteger("5FBFF498...3B7E", 16);
    BigInteger q = new BigInteger("80000000...F5B3", 16);
    Point P = new Point(
        new BigInteger("2", 16),
        new BigInteger("8E2A8A0E...E8FC8", 16)
    );
    DomainParameters params = new DomainParameters(p, a, b, q, P);

    // Entradas estaticas del documento oficial
    BigInteger e = new BigInteger("2DFBC1B3...3EE5", 16); // Hash
    BigInteger d = new BigInteger("7A929ADE...3B28", 16); // Clave privada
    BigInteger k = new BigInteger("77105C9B...EAB3", 16); // Nonce determinista

    // Valores esperados de la firma
    BigInteger expectedR = new BigInteger("41AA28D2...0493", 16);
    BigInteger expectedS = new BigInteger("1456C64B...9C40", 16);

    // 2. Act (Ejecutar)
    Signature firma = signer.sign(e, d, k, params);

    // 3. Assert (Afirmar)
    assertEquals(expectedR, firma.r, "Fallo en R (C = kP)");
    assertEquals(expectedS, firma.s, "Fallo en S (s = (rd + ke) mod q)");
}
\end{lstlisting}

Este enfoque de pruebas permitió validar la exactitud de las implementaciones abstractas, asegurando que cualquier desviación en el cálculo modular o en la multiplicación escalar sobre la curva elíptica fuese detectada de manera inmediata antes de la fase de integración final.

\subsection{Gestión de la Precisión Arbitraria y Convenciones de Memoria}
El estándar GOST exige realizar operaciones algebraicas sobre cuerpos finitos utilizando números de 256 y 512 bits. Dado que los tipos de datos primitivos de Java (como \texttt{long}) están limitados a 64 bits, la arquitectura base se cimentó íntegramente sobre la clase nativa \texttt{java.math.BigInteger}. Esta decisión eliminó el riesgo de desbordamientos aritméticos (\textit{overflows}) y proporcionó rutinas optimizadas para operaciones fundamentales como la multiplicación modular y el cálculo de inversos multiplicativos.

Adicionalmente, se desarrolló una utilidad centralizada (\texttt{GostUtils}) para gobernar el ordenamiento de bytes. Puesto que los estándares criptográficos rusos imponen el formato \textit{Little-Endian} para la representación de vectores en memoria, mientras que la Máquina Virtual de Java (JVM) opera por defecto en \textit{Big-Endian}, esta abstracción centralizada previno inconsistencias sistémicas durante el intercambio de datos entre la función Hash, el Emisor y el Receptor.

\subsection{Gestión de la Precisión Arbitraria y Convenciones de Memoria}
El estándar GOST exige realizar operaciones algebraicas sobre cuerpos finitos utilizando números de 256 y 512 bits. Dado que los tipos de datos primitivos de Java (como \texttt{long}) están limitados a 64 bits, la arquitectura base se cimentó íntegramente sobre la clase nativa \texttt{java.math.BigInteger}. Esta decisión eliminó el riesgo de desbordamientos aritméticos (\textit{overflows}) y proporcionó rutinas optimizadas para operaciones fundamentales como la multiplicación modular y el cálculo de inversos multiplicativos.

Adicionalmente, se desarrolló una utilidad centralizada (\texttt{GostUtils}) para gobernar el ordenamiento de bytes. Puesto que los estándares criptográficos rusos imponen el formato \textit{Little-Endian} para la representación de vectores en memoria, mientras que la Máquina Virtual de Java (JVM) opera por defecto en \textit{Big-Endian}, esta abstracción centralizada previno inconsistencias sistémicas durante el intercambio de datos entre la función Hash, el Emisor y el Receptor.